\documentclass[11pt,a4paper]{article}
\usepackage[utf8]{inputenc}
\usepackage[T1]{fontenc}
\usepackage[margin=2.5cm]{geometry}
\usepackage{booktabs}
\usepackage{amsmath}
\usepackage{hyperref}
\hypersetup{colorlinks=true, linkcolor=blue, urlcolor=blue, citecolor=blue}

\title{\textbf{Net Electricity Demand Forecasting in France}\\
\large Statistical learning for the 2022--2023 sobriety period}
\author{Marcos Lahoz}
\date{\today}

\begin{document}
\maketitle

\begin{abstract}
We forecast the daily \emph{net} electricity demand of France (consumption minus
wind and solar production) over the 2022--2023 ``sobriety period''. A family of
statistical learners --- additive models, an online state-space (Kalman) adaptation,
quantile additive models and tree ensembles --- is trained and combined through online
expert aggregation. Models are evaluated with the pinball loss at quantile $0.8$, the
competition metric, which deliberately penalises under-forecasting. The
Kalman-adapted GAM and the expert aggregation give the most reliable forecasts.
\end{abstract}

\section{Introduction}
Balancing the electricity grid requires accurate short-term demand forecasts; a
mismatch between supply and demand wastes energy or risks blackouts. We study the
\emph{net demand}, the load that renewable (solar and wind) production cannot cover,
\[
\texttt{Net\_demand} = \texttt{Load} - (\texttt{Wind\_power} + \texttt{Solar\_power}),
\]
during the French ``sobriety period'' (September 2022--September 2023), when consumption
patterns shifted. Forecasts are scored with the \textbf{pinball loss at quantile
$\tau=0.8$},
\[
\mathcal{L}_\tau(y,\hat q) = \big(\tau - \mathbf{1}\{y < \hat q\}\big)\,(y-\hat q),
\]
which is asymmetric: under-forecasting the demand is penalised more heavily than
over-forecasting, a conservative choice for grid security.

\section{Data}
Daily observations span 2013 to September 2022 for training, with the sobriety period
held out for prediction. The covariates combine calendar information (time of year,
day of week, public holidays, school breaks, daylight-saving), weather (temperature
and smoothed/min/max variants, wind, nebulosity) and lagged production/consumption
(one-day and one-week lags of load, net demand, solar and wind power). The data
contains no missing values, so all variables are kept.

\section{Methodology}
Each model uses an estimation technique and covariate set suited to it; the final
prediction combines them. While developing the models we validate on the most recent
years before touching the test horizon.

\paragraph{GAM + Kalman filter.} A Generalized Additive Model captures the smooth,
non-linear effects of temperature, time of year and lagged demand. Its term
contributions are then adapted online with a state-space (Kalman) filter whose
variances are tuned by expectation--maximisation, letting the model track the regime
change of the sobriety period. This is the single most reliable model.

\paragraph{Ridge GAM.} A GAM with a double (ridge) penalty and automatic term
selection yields a frugal, robust fit that resists over-fitting from overly flexible
splines.

\paragraph{GAM + Random Forest on residuals.} A Random Forest is fitted to the GAM
residuals (on temperature, time of year, weekday and a time index) to recover
non-linear structure missed by the additive model.

\paragraph{Random Forest.} A tree ensemble over the full covariate set ($\texttt{mtry}$
tuned), strong for capturing interactions but prone to over-smoothing trends.

\paragraph{Quantile GAMs.} Quantile GAMs at $q=0.2$ and $q=0.8$ bound the predictions
and directly target the asymmetric loss.

\paragraph{Expert aggregation.} Finally, all models are combined online with the
\texttt{opera} mixture (MLpol) under the pinball loss, so the ensemble weights adapt
sequentially to whichever expert performs best.

\section{Results}
Table~\ref{tab:results} reports validation performance (RMSE and pinball loss at
$0.8$). The Kalman-adapted GAM gives the best overall trade-off and the lowest RMSE,
while the Random Forest reaches the lowest pinball loss; the online aggregation
combines their strengths for the final submission.

\begin{table}[h]
\centering
\begin{tabular}{lrr}
\toprule
\textbf{Model} & \textbf{RMSE} & \textbf{Pinball@0.8} \\
\midrule
GAM + Kalman              & 1629 & 537 \\
Ridge GAM                 & 1353 & 562 \\
GAM + RF (residuals)      & 1754 & 716 \\
Random Forest             & 2064 & 465 \\
\bottomrule
\end{tabular}
\caption{Validation performance (lower is better).}
\label{tab:results}
\end{table}

\section{Conclusion}
Combining additive models with online Kalman adaptation, quantile regression and tree
ensembles --- and aggregating them online --- produces robust net-demand forecasts
through a period of shifting consumption. The asymmetric pinball objective is handled
explicitly by the quantile GAMs and the aggregation. Natural extensions include richer
state-space dynamics and probabilistic ensembling across quantiles.

\begin{thebibliography}{9}
\bibitem{kaggle} Kaggle. \emph{Net load forecasting --- sobriety period}.
\url{https://www.kaggle.com/competitions/net-load-soberty-period}
\bibitem{mgcv} Wood, S. N. \emph{Generalized Additive Models: An Introduction with R}.
2nd ed., CRC Press, 2017 (\texttt{mgcv}).
\bibitem{qgam} Fasiolo, M. et al. \emph{Fast calibrated additive quantile regression}.
JASA, 2021 (\texttt{qgam}).
\bibitem{opera} Gaillard, P. \& Goude, Y. \emph{opera: Online Prediction by Expert
Aggregation}. R package.
\end{thebibliography}

\end{document}
